\chapter{E05 --- AI Orchestration}

Goal: coordinate each conversation turn safely. Modules: M04, M19. Detailed turn pipeline: see the AI engine diagrams.

\section*{Stories in this epic}
{\small
\begin{longtable}{@{}>{\raggedright\arraybackslash}p{\dimexpr0.090\linewidth-2\tabcolsep\relax}>{\raggedright\arraybackslash}p{\dimexpr0.780\linewidth-2\tabcolsep\relax}>{\raggedright\arraybackslash}p{\dimexpr0.130\linewidth-2\tabcolsep\relax}@{}}
\toprule
\textbf{ID} & \textbf{User story} & \textbf{Sprint} \\
\midrule
\endhead
\bottomrule
\endfoot
US-040 & As a system, I can process a turn through a versioned agent configuration. & S4 \\
US-041 & As a customer, I can ask a follow-up question using conversation context. & S4 \\
US-042 & As a system, I can handle model timeout or provider failure. & S4 \\
US-043 & As an administrator, I can publish and roll back an agent version. & S6 \\
\end{longtable}
}

\Needspace{175pt}
\section{Epic overview}

\mmdfig{epics/E05-ai-orchestration/01-epic-overview}{E05: Epic overview}{fig:epics-e05-ai-orchestration-01-epic-overview-36}


\mmdfigLH{\section{US-040 --- Process a turn through a versioned agent configuration}

\begin{storybox}{User story US-040}
\textbf{Story.} As a system, I can process a turn through a versioned agent configuration.\par
\textbf{Acceptance.} model/config version, latency, and outcome are recorded.
\end{storybox}
}{epics/E05-ai-orchestration/02-us-040-process-a-turn-through-a-versioned-agent-configuratio}{E05: US-040 --- Process a turn through a versioned agent configuration}{fig:cess-a-turn-through-a-versioned-agent-configuratio-37}


\Needspace{640pt}
\section{US-041 --- Follow-up question using conversation context}

\begin{storybox}{User story US-041}
\textbf{Story.} As a customer, I can ask a follow-up question using conversation context.\par
\textbf{Acceptance.} relevant recent context is used without exposing another customer's data.
\end{storybox}

\mmdfig{epics/E05-ai-orchestration/03-us-041-follow-up-question-using-conversation-context}{E05: US-041 --- Follow-up question using conversation context}{fig:-041-follow-up-question-using-conversation-context-38}


\Needspace{640pt}
\section{US-042 --- Handle model timeout or provider failure}

\begin{storybox}{User story US-042}
\textbf{Story.} As a system, I can handle model timeout or provider failure.\par
\textbf{Acceptance.} bounded retry/fallback occurs; no false success message is delivered.
\end{storybox}

\mmdfig{epics/E05-ai-orchestration/04-us-042-handle-model-timeout-or-provider-failure}{E05: US-042 --- Handle model timeout or provider failure}{fig:04-us-042-handle-model-timeout-or-provider-failure-39}


\mmdfigLH{\section{US-043 --- Publish and roll back an agent version}

\begin{storybox}{User story US-043}
\textbf{Story.} As an administrator, I can publish and roll back an agent version.\par
\textbf{Acceptance.} published versions are immutable and audit logged.
\end{storybox}
}{epics/E05-ai-orchestration/05-us-043-publish-and-roll-back-an-agent-version}{E05: US-043 --- Publish and roll back an agent version}{fig:n-05-us-043-publish-and-roll-back-an-agent-version-40}


